% 04-bptt-jacobian.tex - BPTT derived step by step in Jacobian form
%
% Scope: the backward recurrence; each Jacobian; parameter gradients;
% why the same W_hh appears T times in the product.

\section{Lan truyền ngược qua thời gian}
\label{sec:ch01-bptt}

BPTT không phải là một thuật toán mới. Nó là
\tn{lan truyền ngược}{backpropagation} thông thường,
áp dụng lên đồ thị tính toán đã được tháo cuộn của RNN. Điều duy nhất thay
đổi là gradient phải đi qua $\mat{W}_{hh}$ nhiều lần; mỗi lần tại một bước
thời gian; và chính sự lặp lại đó sẽ là nhân vật chính của ba chương tiếp
theo.

\subsection{ gradient theo đầu ra}

Đạo hàm của $L$ theo $\vect{y}_t$ là điểm khởi đầu của mọi thứ. Từ
\eqref{eq:rnn-loss}:

\begin{equation}
  \pd{L}{\vect{y}_t} = \vect{y}_t - \widehat{\vect{y}}_t.
  \label{eq:bptt-dy}
\end{equation}

Đây là một vector $d_y$-chiều. Mỗi \tn{bước thời gian}{timestep} đóng góp
gradient riêng,
không phụ thuộc vào các bước khác; ít nhất là ở tầng này.

\subsection{ gradient theo hidden state}

Từ $\vect{y}_t$ đi xuống $\vect{h}_t$ là một phép biến đổi tuyến tính
$\vect{y}_t = \mat{W}_{hy} \vect{h}_t + \vect{b}_y$. Jacobian của phép này là
chính $\mat{W}_{hy}$:

\[
  \jac{\vect{y}_t}{\vect{h}_t} = \mat{W}_{hy}.
\]

 Gradient truyền từ $\vect{y}_t$ về $\vect{h}_t$ là tích của Jacobian chuyển
vị với gradient đầu ra:

\begin{equation}
  \pd{L}{\vect{h}_t} \quad\text{nhận từ}\quad
  \underbrace{\mat{W}_{hy}\transpose \pd{L}{\vect{y}_t}}_{\text{từ output}}.
  \label{eq:bptt-dh-output}
\end{equation}

Nhưng $\vect{h}_t$ còn một đường gradient nữa: từ tương lai. Vì
$\vect{h}_{t+1}$ phụ thuộc vào $\vect{h}_t$ qua $\mat{W}_{hh}$, gradient
tại bước $t+1$ sẽ truyền ngược về bước $t$. Đây chính là điểm khác biệt duy
nhất giữa BPTT và lan truyền ngược thông thường, và tôi sẽ gọi nó là
\emph{tín hiệu hồi quy}.

\subsection{Tín hiệu hồi quy}

Tại bước $t$, $\vect{h}_t$ ảnh hưởng đến $\vect{a}_{t+1}$ qua
$\mat{W}_{hh} \vect{h}_t$, rồi $\vect{a}_{t+1}$ qua $\tanh$ thành
$\vect{h}_{t+1}$. Ta cần Jacobian của $\vect{h}_{t+1}$ theo $\vect{h}_t$.

Trước hết, đạo hàm của $\tanh$ là một ma trận đường chéo:

\[
  \jac{\vect{h}_{t+1}}{\vect{a}_{t+1}}
  = \diag\bigl(1 - \vect{h}_{t+1}^2\bigr),
\]

với $\vect{h}_{t+1}^2$ là bình phương từng phần tử. Ký hiệu $\vect{1}$ là
vector toàn 1, và phép bình phương được áp dụng element-wise.

Tiếp theo, $\jac{\vect{a}_{t+1}}{\vect{h}_t} = \mat{W}_{hh}$, vì
$\vect{a}_{t+1} = \mat{W}_{xh} \vect{x}_{t+1} + \mat{W}_{hh} \vect{h}_t
+ \vect{b}_h$.

Theo quy tắc dây chuyền Jacobian:

\begin{equation}
  \jac{\vect{h}_{t+1}}{\vect{h}_t}
  = \jac{\vect{h}_{t+1}}{\vect{a}_{t+1}}
    \;\jac{\vect{a}_{t+1}}{\vect{h}_t}
  = \diag\bigl(1 - \vect{h}_{t+1}^2\bigr) \; \mat{W}_{hh}.
  \label{eq:bptt-dhnext-dht}
\end{equation}

Khi truyền gradient, ta nhân với chuyển vị:

\[
  \Bigl(\jac{\vect{h}_{t+1}}{\vect{h}_t}\Bigr)\transpose
  = \mat{W}_{hh}\transpose \;
    \diag\bigl(1 - \vect{h}_{t+1}^2\bigr).
\]

\subsection{Công thức truy hồi cho hidden state}

Gộp hai đường gradient; từ output và từ tương lai; ta có công thức truy hồi
trung tâm của BPTT:

\begin{equation}
  \boxed{
    \pd{L}{\vect{h}_t}
    = \mat{W}_{hy}\transpose \pd{L}{\vect{y}_t}
    + \mat{W}_{hh}\transpose \;
      \Bigl(
        \pd{L}{\vect{h}_{t+1}}
        \odot \bigl(\vect{1} - \vect{h}_{t+1}^2\bigr)
      \Bigr)
  }.
  \label{eq:bptt-dh-recurrence}
\end{equation}

Trong đó $\odot$ là tích Hadamard (nhân từng phần tử). Tôi viết hạng tử thứ
hai dưới dạng Hadamard thay vì $\diag$ để gần với code hơn: trong Python,
\mintinline{python}{dh_next * (1 - h_next**2)} rồi nhân với
\mintinline{python}{W_hh.T} là ba dòng, không cần dựng ma trận đường chéo
$d_h \times d_h$.

Điều kiện biên: tại bước cuối cùng $t = T$, không có tương lai để nhận
gradient:

\begin{equation}
  \pd{L}{\vect{h}_T}
  = \mat{W}_{hy}\transpose \pd{L}{\vect{y}_T}.
  \label{eq:bptt-dhT}
\end{equation}

Từ đây, ta chạy \eqref{eq:bptt-dh-recurrence} lùi dần: $t = T-1, T-2, \dots,
1$. Ở mỗi bước, gradient $\pd{L}{\vect{h}_t}$ nhận thêm tín hiệu từ output
hiện tại và từ tất cả các bước tương lai đã được gộp qua $\mat{W}_{hh}$.

\subsection{ gradient qua $\tanh$}

Một khi đã có $\pd{L}{\vect{h}_t}$, bước qua $\tanh$ là đạo hàm element-wise:

\begin{equation}
  \pd{L}{\vect{a}_t}
  = \pd{L}{\vect{h}_t}
    \odot \bigl(\vect{1} - \vect{h}_t^2\bigr).
  \label{eq:bptt-da}
\end{equation}

Đây là vector $d_h$-chiều. Nó là ``tín hiệu lỗi'' tại pre-activation của
bước $t$, và sẽ được dùng để tính gradient cho mọi tham số đóng góp vào
$\vect{a}_t$.

\subsection{ gradient cho tham số}

Từ $\pd{L}{\vect{a}_t}$ và $\pd{L}{\vect{y}_t}$, gradient cho năm tham số
được tích lũy qua toàn bộ \tn{chuỗi}{sequence}:

\begin{align}
  \pd{L}{\mat{W}_{xh}} &= \sum_{t=1}^{T}
    \pd{L}{\vect{a}_t} \otimes \vect{x}_t\transpose,
    \label{eq:bptt-dWxh} \\
  \pd{L}{\mat{W}_{hh}} &= \sum_{t=1}^{T}
    \pd{L}{\vect{a}_t} \otimes \vect{h}_{t-1}\transpose,
    \label{eq:bptt-dWhh} \\
  \pd{L}{\mat{W}_{hy}} &= \sum_{t=1}^{T}
    \pd{L}{\vect{y}_t} \otimes \vect{h}_t\transpose,
    \label{eq:bptt-dWhy} \\
  \pd{L}{\vect{b}_h} &= \sum_{t=1}^{T}
    \pd{L}{\vect{a}_t},
    \label{eq:bptt-dbh} \\
  \pd{L}{\vect{b}_y} &= \sum_{t=1}^{T}
    \pd{L}{\vect{y}_t}.
    \label{eq:bptt-dby}
\end{align}

Ký hiệu $\vect{u} \otimes \vect{v}\transpose$ là tích ngoài: nếu $\vect{u}
\in \R^m$ và $\vect{v} \in \R^n$, kết quả là ma trận $m \times n$ với phần
tử $(i,j) = u_i v_j$. Đây chính là dạng Jacobian của gradient theo
\tn{ma trận trọng số}{weight matrix}: mỗi cặp (đầu vào, tín hiệu lỗi) đóng góp
một hạng tử vào tổng.

\subsection{Vì sao $\mat{W}_{hh}$ lặp $T$ lần}

Nhìn vào \eqref{eq:bptt-dh-recurrence}: $\mat{W}_{hh}\transpose$ xuất hiện ở
mọi bước lùi. Nếu bạn tháo cuộn công thức này từ $t = T$ về $t = 1$, bạn sẽ
thấy $\pd{L}{\vect{h}_1}$ chứa $\mat{W}_{hh}\transpose$ nhân với chính nó $T$
lần. Khi $T$ lớn, tích này hoặc co về không hoặc bùng nổ, tùy vào
\tn{trị riêng}{eigenvalue} lớn nhất của $\mat{W}_{hh}$. Đó chính là vấn đề
\tn{triệt tiêu gradient}{vanishing gradient} và
\tn{bùng nổ gradient}{exploding gradient}; chủ đề của chương~2 và~3.

Tôi sẽ không dẫn xuất điều kiện đó ở đây, vì chương~3 dành trọn một chương
cho nó. Điều cần nhớ lúc này: BPTT là một vòng lặp, và mỗi vòng lặp nhân thêm
một ma trận $\mat{W}_{hh}$. Đó là toàn bộ câu chuyện.

\subsection{Tóm tắt thuật toán}

Toàn bộ BPTT gói gọn trong năm bước:

\begin{enumerate}
  \item \textbf{Lan truyền xuôi}: tính và lưu $\vect{x}_t$, $\vect{a}_t$,
    $\vect{h}_t$, $\vect{y}_t$ cho $t = 1, \dots, T$.
  \item \textbf{Khởi tạo}: $\pd{L}{\vect{h}_{T+1}} = \vect{0}$ (không có
    tương lai sau bước cuối).
  \item \textbf{Lùi từ $T$ về $1$}:
    \begin{itemize}
      \item $\pd{L}{\vect{y}_t} = \vect{y}_t - \widehat{\vect{y}}_t$
      \item $\pd{L}{\vect{h}_t}$ từ \eqref{eq:bptt-dh-recurrence}
      \item $\pd{L}{\vect{a}_t}$ từ \eqref{eq:bptt-da}
      \item Tích lũy gradient tham số từ
        \eqref{eq:bptt-dWxh} đến \eqref{eq:bptt-dby}.
    \end{itemize}
  \item \textbf{Kết thúc}: tổng gradient cho năm tham số đã sẵn sàng cho
    bước cập nhật (SGD, Adam, v.v.).
\end{enumerate}

Trong code, bước~3 là một vòng lặp \mintinline{python}{for t in reversed(range(T))}.
Mọi thứ khác là phép nhân ma trận và Hadamard product. Toàn bộ cài đặt dài
chưa đến 30 dòng Python; bạn sẽ thấy nó trong mục tiếp theo.
